\documentclass[UTF8,fontset=fandol,11pt]{ctexart}

\usepackage[a4paper,margin=2.5cm]{geometry}
\usepackage{amsmath,amssymb,amsthm}
\usepackage{hyperref}
\hypersetup{
  unicode=true,
  colorlinks=true,
  linkcolor=blue,
  pdftitle={Motzkin交错和的根序列严格对数凹性},
  pdfauthor={}
}

\newtheorem{theorem}{定理}
\newtheorem{lemma}{引理}
\newcommand{\E}{\mathbb{E}}
\DeclareMathOperator{\Var}{Var}
\DeclareMathOperator{\Cov}{Cov}

\title{Motzkin 交错和的根序列严格对数凹性}
\author{}
\date{}

\begin{document}
\maketitle

\begin{abstract}
本文从定义递推构造 Motzkin 数的积分表示，将其交错和分解为一个正积分主项与绝对值小于 $1$ 的误差。通过与 Gamma 分布比较，证明主项根对数的连续插值具有明确的二阶导数负上界，再证明该凹性余量严格超过交错误差的二阶差分。由此证明从第 $9$ 项起的根序列严格对数凹。全部阈值和常数均用显式不等式验证。
\end{abstract}

\begin{theorem}\label{thm:main}
设数列 $(M_n)_{n\ge0}$ 由
\[
M_0=M_1=1,
\qquad
(n+2)M_n=(2n+1)M_{n-1}+3(n-1)M_{n-2}
\quad(n\ge2)
\]
定义，并令
\[
T_n=\sum_{k=0}^{n}(-1)^{n-k}M_k.
\]
则 $T_n>0$ 对所有整数 $n\ge9$ 成立，且
\begin{equation}\label{eq:conclusion}
\left(T_n^{1/n}\right)^2
>T_{n-1}^{1/(n-1)}T_{n+1}^{1/(n+1)}
\qquad(n\ge10).
\end{equation}
因此 $(T_n^{1/n})_{n\ge9}$ 是严格对数凹序列。
\end{theorem}

全文中的 $\log$ 均表示自然对数。我们使用标准的 Beta 函数与 Gamma 函数
\[
B(u,v)=\int_0^1 y^{u-1}(1-y)^{v-1}\,dy,
\qquad
\Gamma(u)=\int_0^\infty s^{u-1}e^{-s}\,ds
\quad(u,v>0),
\]
以及恒等式
\[
B(u,v)=\frac{\Gamma(u)\Gamma(v)}{\Gamma(u+v)},
\qquad
\Gamma\left(\frac12\right)=\sqrt\pi,
\qquad
\Gamma\left(\frac32\right)=\frac{\sqrt\pi}{2}.
\]

\section{积分表示及有界误差}

定义
\[
q(x)=(3-x)(x+1)=3+2x-x^2,
\qquad
I_n=\frac1{2\pi}\int_{-1}^{3}x^n\sqrt{q(x)}\,dx
\quad(n\ge0).
\]
半径为 $2$ 的半圆面积给出 $I_0=1$。由于 $\sqrt{q(x)}$ 关于 $x=1$ 对称，有 $I_1=I_0=1$。对 $n\ge2$，函数 $x^{n-1}q(x)^{3/2}$ 在两个积分端点均为零。因此
\begin{align*}
0
&=\frac1{2\pi}\int_{-1}^{3}
  \frac{d}{dx}\bigl(x^{n-1}q(x)^{3/2}\bigr)\,dx\\
&=(n-1)\bigl(3I_{n-2}+2I_{n-1}-I_n\bigr)
  +3(I_{n-1}-I_n)\\
&=3(n-1)I_{n-2}+(2n+1)I_{n-1}-(n+2)I_n.
\end{align*}
故 $I_n$ 与 $M_n$ 满足相同初值和递推，从而
\begin{equation}\label{eq:moment}
M_n=\frac1{2\pi}\int_{-1}^{3}x^n\sqrt{(3-x)(x+1)}\,dx.
\end{equation}

有限几何级数给出
\[
\sum_{k=0}^{n}(-1)^{n-k}x^k
=\frac{x^{n+1}+(-1)^n}{x+1}
\quad(x\ne-1).
\]
另一方面，作代换 $x+1=4u$ 可得
\[
\frac1{2\pi}\int_{-1}^{3}
\sqrt{\frac{3-x}{x+1}}\,dx
=\frac2\pi B\left(\frac12,\frac32\right)=1.
\]
将此式及有限几何级数代入 \eqref{eq:moment}，得到
\begin{equation}\label{eq:T-integral}
T_n=(-1)^n+\frac1{2\pi}\int_{-1}^{3}
x^{n+1}\sqrt{\frac{3-x}{x+1}}\,dx.
\end{equation}
这里端点 $x=-1$ 附近的奇性为 $O((x+1)^{-1/2})$，可积，因而拆分积分合法。

对实数 $t>-2$，定义正函数
\begin{equation}\label{eq:P}
P(t)=\frac1{2\pi}\int_0^3
x^{t+1}\sqrt{\frac{3-x}{x+1}}\,dx.
\end{equation}
对整数 $n\ge0$，再定义
\[
A_n=\frac1{2\pi}\int_0^1
u^{n+1}\sqrt{\frac{3+u}{1-u}}\,du.
\]
在 \eqref{eq:T-integral} 的负半轴部分作代换 $x=-u$，便有
\begin{equation}\label{eq:decomposition}
T_n=P(n)+(-1)^n(1-A_n).
\end{equation}
利用 $\sqrt{3+u}\le2$ 及 $u^{n+1}\le u$，得到
\[
0<A_n
\le\frac1\pi B\left(n+2,\frac12\right)
\le\frac1\pi B\left(2,\frac12\right)
=\frac4{3\pi}<1.
\]
最后一步使用 $\pi>3$。因此
\begin{equation}\label{eq:error-one}
|T_n-P(n)|<1
\qquad(n\ge0).
\end{equation}

\section{与 Gamma 分布比较}

先记录一个简单的概率比较。这里 $\E$、$\Var$ 和 $\Cov$ 分别表示数学期望、方差和协方差；下标标明所用的概率密度。

\begin{lemma}\label{lem:comparison}
设 $g$ 是概率密度，$r$ 是正的递减函数，并且 $0<\E_g r<\infty$。将 $g(s)r(s)$ 归一化，得到概率密度
\[
p(s)=\frac{g(s)r(s)}{\E_g r}.
\]
若 $\varphi$ 为递增函数，且下述期望均有限，则
\[
\E_p\varphi\le\E_g\varphi.
\]
\end{lemma}

\begin{proof}
取相互独立且密度均为 $g$ 的随机变量 $X,Y$。由于 $\varphi$ 递增、$r$ 递减，
\[
2\Cov_g(\varphi,r)
=\E\bigl[(\varphi(X)-\varphi(Y))(r(X)-r(Y))\bigr]\le0.
\]
因此 $\E_g(\varphi r)\le\E_g\varphi\,\E_g r$，除以 $\E_g r$ 即得结论。
\end{proof}

本节及下一节均设实数 $t\ge9$。令
\begin{equation}\label{eq:ZaC}
Z(t)=\frac{P(t)}{3^t},
\qquad
a=t+\frac{13}{8},
\qquad
C=\frac{3^{5/2}}{8\sqrt\pi}.
\end{equation}
此时总有 $a>0$。在 \eqref{eq:P} 中先作 $x=3y$，再作 $y=e^{-s}$，得到
\begin{align}
Z(t)
&=\frac{3^{5/2}}{2\pi}\int_0^1
y^{t+1}\frac{\sqrt{1-y}}{\sqrt{1+3y}}\,dy\notag\\
&=\frac{3^{5/2}}{2\pi}\int_0^\infty
e^{-(t+2)s}\frac{\sqrt{1-e^{-s}}}{\sqrt{1+3e^{-s}}}\,ds\notag\\
&=\frac{3^{5/2}}{4\pi}\int_0^\infty
e^{-as}s^{1/2}r(s)\,ds, \label{eq:Z-Gamma}
\end{align}
其中
\begin{equation}\label{eq:r}
r(s)=\sqrt{\frac{1-e^{-s}}s}\,
\frac{2e^{-3s/8}}{\sqrt{1+3e^{-s}}}
\qquad(s>0).
\end{equation}
显然 $r(s)>0$，且 $r(0+)=1$。直接求导可得
\begin{equation}\label{eq:r-derivative}
\frac{r'(s)}{r(s)}
=\frac12\left(\frac1{e^s-1}-\frac1s\right)
-\frac38+\frac{3e^{-s}}{2(1+3e^{-s})}<0.
\end{equation}
事实上，$e^s>1+s$ 使括号内严格为负；令 $z=e^{-s}\le1$，则
\[
\frac{3z}{2(1+3z)}\le\frac38
\quad\Longleftrightarrow\quad z\le1,
\]
所以最后两项之和不大于零。由此得到 $r$ 递减且 $0<r(s)\le1$。

令 $g_a$ 为形状参数 $3/2$、速率参数 $a$ 的 Gamma 概率密度，即
\begin{equation}\label{eq:gamma-density}
g_a(s)=\frac{a^{3/2}}{\Gamma(3/2)}s^{1/2}e^{-as}
\quad(s>0),
\qquad
\int_0^\infty g_a(s)\,ds=1.
\end{equation}
再令 $S$ 具有概率密度
\begin{equation}\label{eq:S-density}
p_t(s)=\frac{g_a(s)r(s)}{\displaystyle\int_0^\infty g_a(v)r(v)\,dv}
=\frac{e^{-as}s^{1/2}r(s)}{\displaystyle\int_0^\infty e^{-av}v^{1/2}r(v)\,dv}.
\end{equation}
利用 $r\le1$，由 \eqref{eq:Z-Gamma} 得到
\begin{equation}\label{eq:Z-upper}
Z(t)\le\frac{3^{5/2}}{4\pi}\frac{\Gamma(3/2)}{a^{3/2}}
=C a^{-3/2}.
\end{equation}
将引理 \ref{lem:comparison} 分别用于递增函数 $s$ 与 $s^2$，并使用 Gamma 分布的矩，得到
\begin{equation}\label{eq:S-moments}
\E S\le\frac3{2a},
\qquad
\E S^2\le\frac{15}{4a^2}.
\end{equation}
这两个 Gamma 矩也可直接将 \eqref{eq:gamma-density} 分别乘以 $s$ 和 $s^2$ 后积分求出。

\section{连续主项的凹性余量}

对于任意 $t_0\ge9$，可在 $t_0$ 的一个紧邻域内取正常数 $c$，使 $a=t+13/8\ge c$。由于 $0<r\le1$，\eqref{eq:Z-Gamma} 的被积函数及关于 $t$ 的前两阶导数，在无穷远处分别由常数乘以
\[
s^{1/2}e^{-cs},\quad s^{3/2}e^{-cs},\quad s^{5/2}e^{-cs}
\]
控制；在零点附近这些函数也可积。因此可在积分号下求两次导数，并由 \eqref{eq:S-density} 得
\begin{equation}\label{eq:logZ-derivatives}
(\log Z)'(t)=-\E S,
\qquad
(\log Z)''(t)=\Var S.
\end{equation}

定义
\[
f(t)=\frac{\log P(t)}t
=\log3+\frac{\log Z(t)}t
\qquad(t\ge9).
\]
由 \eqref{eq:logZ-derivatives}、\eqref{eq:Z-upper}、\eqref{eq:S-moments} 及 $\Var S\le\E S^2$，有
\begin{align}
t^3f''(t)
&=t^2\Var S+2t\E S+2\log Z(t)\notag\\
&\le R(t), \label{eq:fpp-R}
\end{align}
其中
\begin{equation}\label{eq:R}
R(t)=\frac{15}{4}\left(\frac ta\right)^2
+3\frac ta+2\log C-3\log a.
\end{equation}

下面用显式常数证明 $R(t)<-3/2$。记 $c_0=13/8$，则 $a=t+c_0$。对 $t\ge9$，有 $0<t/a<1$，从而
\begin{align}
R'(t)
&=\left(\frac{15}{2}\frac ta+3\right)\frac{c_0}{a^2}-\frac3a\notag\\
&\le\frac{(21/2)c_0-3a}{a^2}<0. \label{eq:R-decreasing}
\end{align}
最后一步是因为 $a\ge85/8$，而
\[
\frac{21}{2}c_0=\frac{273}{16}
<\frac{255}{8}=3\cdot\frac{85}{8}.
\]

我们只需验证 $t=9$。所用初等常数估计为
\begin{equation}\label{eq:elementary-constants}
3<\pi<\frac72,
\qquad
e<\frac{11}{4}.
\end{equation}
其中 $\pi>3$ 来自单位圆内接正六边形的周长严格小于圆周长；上界可由
\[
\pi=4\int_0^1\frac{dx}{1+x^2}
<4\int_0^1(1-x^2+x^4)\,dx
=\frac{52}{15}<\frac72
\]
得到。这里 $(1+x^2)^{-1}<1-x^2+x^4$ 对 $x>0$ 成立。对于 $e$，由 $(4+j)!\ge4!\,5^j$ 有
\[
e=\sum_{k=0}^\infty\frac1{k!}
\le\frac83+\frac1{24}\sum_{j=0}^\infty5^{-j}
=\frac{87}{32}<\frac{11}{4}.
\]
由 $\pi>3$ 及 $\log(1+x)<x$（$x>0$），得到
\begin{equation}\label{eq:C-log}
C<\frac98,
\qquad 2\log C<\frac14.
\end{equation}
又因为
\[
11^7=19\,487\,171
<19\,652\,000=32\cdot85^3,
\]
所以
\[
e^7<\left(\frac{11}{4}\right)^7
<\left(\frac{85}{8}\right)^3,
\qquad
\log\left(\frac{85}{8}\right)>\frac73.
\]
在 $t=9$ 时，$a=85/8$，$t/a=72/85$，从而
\begin{align*}
R(9)
&<\frac{15}{4}\left(\frac{72}{85}\right)^2
+\frac{216}{85}+\frac14-7\\
&=\frac{1512}{289}-\frac{27}{4}
=-\frac{1755}{1156}<-\frac32.
\end{align*}
结合 \eqref{eq:R-decreasing} 和 \eqref{eq:fpp-R}，得到关键估计
\begin{equation}\label{eq:fpp-negative}
\boxed{\displaystyle
f''(t)<-\frac3{2t^3}\qquad(t\ge9).}
\end{equation}

\section{主项的统一下界}

当 $0\le y\le1$ 时，$(1+3y)^{-1/2}\ge1/2$。因此，对整数 $m\ge9$，由 \eqref{eq:P} 中的代换 $x=3y$ 得
\begin{align}
P(m)
&\ge\frac{3^{m+5/2}}{4\pi}B\left(m+2,\frac32\right)\notag\\
&=C3^m\frac{\Gamma(m+2)}{\Gamma(m+7/2)}. \label{eq:P-beta-lower}
\end{align}

为控制 Gamma 函数之比，令 $b>0$，并令随机变量 $X$ 具有形状参数 $b$、速率参数 $1$ 的 Gamma 密度
\[
\frac{x^{b-1}e^{-x}}{\Gamma(b)}\qquad(x>0).
\]
于是 $\E X^2=b(b+1)$。将 Jensen 不等式用于凹函数 $v^{3/4}$，得到
\[
\frac{\Gamma(b+3/2)}{\Gamma(b)}
=\E X^{3/2}
\le(\E X^2)^{3/4}
=[b(b+1)]^{3/4}<(b+1)^{3/2}.
\]
取 $b=m+2$。由 \eqref{eq:elementary-constants}，
\[
C^2=\frac{243}{64\pi}>\frac{243}{224}>1,
\]
所以 $C>1$，进而 \eqref{eq:P-beta-lower} 给出
\begin{equation}\label{eq:P-lower}
\boxed{\displaystyle
P(m)>L_m:=\frac{3^m}{(m+3)^{3/2}}\qquad(m\ge9).}
\end{equation}

\section{离散误差控制与结论}

对于整数 $n\ge10$，定义
\[
U_n=\frac{3^{n-1}}{(n+2)^{3/2}}=L_{n-1},
\qquad
V_n=1+\frac83\frac{(n+1)^3}{n-1}.
\]
我们先证明
\begin{equation}\label{eq:U-V}
U_n>V_n\qquad(n\ge10).
\end{equation}
基础步用 $\sqrt{12}<4$ 得到
\[
U_{10}>\frac{3^9}{48}=\frac{6561}{16}
>\frac{10675}{27}=V_{10}.
\]
中间的有理不等式等价于 $177147>170800$。对 $n\ge10$，由于 $0<(n+2)/(n+3)<1$，
\begin{equation}\label{eq:U-growth}
\frac{U_{n+1}}{U_n}
=3\left(\frac{n+2}{n+3}\right)^{3/2}
>3\left(\frac{n+2}{n+3}\right)^2
\ge3\left(\frac{12}{13}\right)^2
=\frac{432}{169}>2.
\end{equation}
另一方面，记 $D_n=V_n-1$，则
\[
\frac{D_{n+1}}{D_n}
=\left(\frac{n+2}{n+1}\right)^3\frac{n-1}{n}
<\left(\frac{12}{11}\right)^3<2,
\]
其中最后一步为 $1728<2662$。所以
\[
V_{n+1}=1+D_{n+1}<1+2D_n<2V_n.
\]
结合 \eqref{eq:U-growth}，归纳得 \eqref{eq:U-V}。同一比值估计还说明 $L_m$ 在整数 $m\ge9$ 上递增。

由 \eqref{eq:P-lower} 和 \eqref{eq:U-V}，有 $P(j)>1$ 对所有 $j\ge9$ 成立。再由 \eqref{eq:error-one}，得到 $T_j>0$。故以下对数均有定义。令
\[
\delta_j=\frac{\log T_j}{j}-f(j)
\qquad(j\ge9).
\]
对 $|u|<1$，由 $\log(1+u)=\int_0^u(1+v)^{-1}\,dv$ 得
\[
|\log(1+u)|\le\frac{|u|}{1-|u|}.
\]
取 $u=(T_j-P(j))/P(j)$，并使用 \eqref{eq:error-one}，得
\begin{equation}\label{eq:delta}
|\delta_j|<\frac1{j(P(j)-1)}.
\end{equation}
固定 $n\ge10$。当 $j\in\{n-1,n,n+1\}$ 时，由 \eqref{eq:P-lower} 及 $L_j\ge L_{n-1}=U_n$，得到
\[
|\delta_j|<\frac1{(n-1)(U_n-1)}.
\]
所以由 \eqref{eq:U-V}，
\begin{align}
|2\delta_n-\delta_{n-1}-\delta_{n+1}|
&<\frac4{(n-1)(U_n-1)}\notag\\
&<\frac3{2(n+1)^3}. \label{eq:error-second}
\end{align}

另一方面，两次微积分基本定理给出二阶差分的积分恒等式
\[
f(n-1)-2f(n)+f(n+1)
=\int_{-1}^{1}(1-|s|)f''(n+s)\,ds.
\]
由于 $n+s\ge9$，可用 \eqref{eq:fpp-negative}，于是
\begin{align}
2f(n)-f(n-1)-f(n+1)
&>\frac32\int_{-1}^{1}\frac{1-|s|}{(n+s)^3}\,ds\notag\\
&\ge\frac3{2(n+1)^3}. \label{eq:main-second}
\end{align}
最后一步使用 $n+s\le n+1$ 以及
$\int_{-1}^{1}(1-|s|)\,ds=1$。

现在将 \eqref{eq:error-second} 与 \eqref{eq:main-second} 合并，得到
\begin{align*}
\frac{2\log T_n}{n}
-\frac{\log T_{n-1}}{n-1}
-\frac{\log T_{n+1}}{n+1}
&=2f(n)-f(n-1)-f(n+1)\\
&\quad+2\delta_n-\delta_{n-1}-\delta_{n+1}\\
&>0
\qquad(n\ge10).
\end{align*}
指数化即得 \eqref{eq:conclusion}，定理 \ref{thm:main} 得证。
\hfill$\square$

\end{document}
